\chapter{质量、安全与可靠性设计}

\section{非功能指标}

赛题要求\textbf{事实问答准确率不低于90\%}、\textbf{语音问答时延合理}（示例阈值小于5秒）且系统稳定。设计上将指标拆分为可脚本验证的“准确、检索、构建、接口”与需真实浏览器和云账号验证的“口型、表情、云端时延”，避免以\textbf{离线Mock替代现场自然度结论}。

\begin{table}[H]
  \centering
  \caption{非功能指标与实现机制表}\label{tab:nfr}
  \begin{tabularx}{\textwidth}{p{0.16\textwidth}p{0.20\textwidth}X p{0.21\textwidth}}
    \toprule
    质量属性 & 目标 & 设计机制 & 验证方式 \\
    \midrule
    准确性 & 事实问答准确率$\geq90\%$ & 本地知识边界、混合检索、引用、无资料拒答 & 50题标准问答集 \\
    检索能力 & Top-3覆盖正确证据 & 实体、词项、意图、覆盖度加权与稳定重排 & 20条检索用例 \\
    响应性 & 语音链路目标$<5$秒 & FAQ缓存、分段TTS、阶段耗时、超时重试 & 现场云端实测与接口指标 \\
    稳定性 & 无崩溃、无长时间无响应 & 全局错误边界、请求日志、重复提交锁、演示降级 & 并发与降级脚本 \\
    可用性 & 移动端核心流程完整 & 五Tab信息架构、安全区、加载态、可恢复错误提示 & APK、小程序与HAP人工走查 \\
    跨平台性 & 同一业务覆盖Web、桌面和游客多端 & API合同复用，录音、定位、地图与WebView按端适配 & 各目标编译、包体校验与代表性设备运行 \\
    可维护性 & 内容和供应商可替换 & 服务模块、OpenAPI、知识重建、LLM/ASR/TTS适配器 & 工程检查与生产构建 \\
    \bottomrule
  \end{tabularx}
\end{table}

\section{安全与可信设计}

AI系统风险管理需要在治理、场景映射、测量和处置之间形成\textbf{持续闭环}\upcite{nist2023airmf}。本项目据此设置\textbf{四类控制}：输入侧敏感与越界检测，知识侧来源约束，生成侧引用和拒答，运营侧人工标注与知识修正。系统提示词明确\textbf{“仅基于灵山资料作答”}，不会用未检索到的模型记忆补全票价、开放时间等事实。

管理端使用\textbf{Bearer Token}\textbf{隔离游客权限}；密码采用\textbf{随机盐和scrypt派生摘要}。scrypt通过内存难解设计提高定制硬件暴力尝试的成本\upcite{percival2009scrypt}。云端凭证保存在环境变量，\textbf{不进入Git}；日志避免记录密钥原文，数据库备份与恢复脚本用于演示前保护数据状态。

\section{异常与降级策略}

\begin{table}[H]
  \centering
  \caption{关键风险与降级应对策略表}\label{tab:risk}
  \begin{tabularx}{\textwidth}{p{0.20\textwidth}p{0.18\textwidth}X p{0.20\textwidth}}
    \toprule
    风险事件 & 影响 & 系统应对 & 用户侧结果 \\
    \midrule
    讯飞形象凭证、额度或网络异常 & 云端视频流、口型或语音不可用 & 捕获SDK错误并释放实例；运营人员可启用\textbf{Live2D本地模型}，继续异常则切换文本/音频模式 & 角色切换过程可控，问答与路线仍可使用 \\
    大模型超时或密钥缺失 & 无法生成自然语言答案 & 超时重试；演示模式调用\textbf{本地确定性提供者} & 获得稳定的有据答案 \\
    ASR不可用或浏览器不支持 & 无法完成语音输入 & 显示文字输入，并支持模拟转写演示 & 用户可继续文本提问 \\
    TTS长文本或失败 & 首次播放慢、播报中断 & 断句分段合成，先返回可播放片段 & 答案文本立即可读 \\
    腾讯地图离线、Key或SDK异常 & 在线底图无法加载 & Web恢复Leaflet/手绘层；游客端保留原生地图与本地手绘图切换 & 点位、路线、搜索和讲解仍可使用 \\
    定位权限拒绝或GPS漂移 & 无法自动居中或最近点判断偏差 & 保留当前中心、搜索与手动选点；最近点只作为建议 & 游客可手工选择景点继续导览 \\
    移动端WebView或平台能力受限 & 云数字人或在线地图在特定端不可用 & 微信使用本地数字人表现；HarmonyOS等端保留手绘地图和TTS & 核心问答、地图点位与路线不中断 \\
    知识库未构建或低相关 & 可能无证据或产生幻觉 & 返回未构建状态；低相关问题进入\textbf{诚实拒答} & 不编造景区事实 \\
    数据库或演示数据损坏 & 历史、配置和后台不可用 & 备份恢复；必要时重建演示数据库 & 可恢复到预置状态 \\
    \bottomrule
  \end{tabularx}
\end{table}

\section{可观测与运维}

每次HTTP请求记录方法、路径、状态与耗时；问答响应额外记录\textbf{检索和LLM耗时}，语音响应记录ASR、RAG、LLM、TTS及总耗时。管理员可从会话详情查看\textbf{答案标签、情绪、时延与反馈}，并对错误消息标注“错误”或“需补知识”。\filepath{env:check}、\filepath{backup}、\filepath{restore}、\filepath{test:concurrency}和\filepath{test:demo-mode}脚本共同覆盖\textbf{部署前检查、数据保护和现场演示保障}。桌面交付另执行Web构建、ASAR树、\filepath{node:sqlite}运行时、系统架构和安装包后缀校验；移动交付记录AppID、包标识、权限声明、版本号和产物路径。

\section{隐私与数据最小化}

游客侧以\textbf{匿名会话标识}工作，不要求采集身份证、手机号等身份数据；定位只在游客主动点击定位后请求，录音只在按下录音按钮并授予麦克风权限后上传，地图与语音服务不把原始坐标、音频或云端凭证写入前端持久化配置。运营大屏处理年龄段、性别、消费和满意度等统计字段，\textbf{输出聚合趋势而非个人画像明细}。正式落地时应增加明确告知、保存期限、日志脱敏、角色最小权限和备份加密策略，并为语音上传与云端数字人处理提供\textbf{用户授权入口}。
